\begin{textsample}{POS Dim 1 – vem\_ed\_18 – Score 5.00 – vem\_ed\_18/t084.txt}  \label{ex:f1_pos_115}
UM FAROL DE ESPERANÇA CRÍTICA Hope Windle, diretora do SUNY COIL Center (State University of New York / EUA), \textbf{referência} mundial em Intercâmbios Virtuais, encerrou o ciclo de apresentações.
Após a exibição de um vídeo com trechos da entrevista concedida por videoconferência para VEm 17 em 23 de março de 2023, Hope Windle respondeu a perguntas da audiência.
No vídeo gravado, Hope \textbf{mencionou} \textbf{que} mais \textbf{pessoas} do setor de design instrucional, de \textbf{apoio} pedagógico e da área de ensino têm vindo participar de projetos COIL na SUNY, enquanto no \textbf{início} da pandemia \textbf{eram} executivos internacionais e \textbf{pessoas} ligadas à internacionalização no currículo.
Sobre o \textbf{que} vislumbra para os próximos 5 anos, ela \textbf{é} entusiasmada com o \textbf{que} virá. “Nos anos 2000 a 2010, as \textbf{pessoas} \textbf{que} primeiro adotaram esses projetos se perguntavam: como podemos usar essas tecnologias?
Agora, as \textbf{pessoas} estão \textbf{reconhecendo} COIL, \textbf{que} não \textbf{é} algo para as \textbf{pessoas} sem condições econômicas \textbf{que} não podem viajar e sim como um farol de esperança \textbf{crítica} para \textbf{que} as \textbf{pessoas} possam repensar sua educação e o \textbf{que} vão fazer com suas vidas”.
Para contextualizar esse conceito em voga na contemporaneidade, fez a seguinte citação:"A esperança \textbf{crítica} reflete a \textbf{habilidade} de avaliar o \textbf{ambiente} de maneira realista, através de uma lente de equidade e justiça, ao mesmo tempo em \textbf{que} vislumbra a possibilidade de um \textbf{futuro} melhor"(DUGAN, 2017; DUNCAN- ANDRADE, 2009).
Hope \textbf{mencionou} ainda a \textbf{questão} dolorosa da escravidão \textbf{que} afeta discussões sobre classe e raça, tanto nos Estados Unidos \textbf{quanto} no Brasil.
Pontuou questões-chave para a educação global na atualidade e no \textbf{futuro} próximo, como o pensamento inclusivo e a interseccionalidade.
A diretora do SUNY COIL Center citou Abby Bryant, diretora do \textbf{programa} EOP, grupo de \textbf{apoio} acadêmico, psicológico e financeiro para \textbf{pessoas} com dificuldades socioeconômicas em universidades dos EUA.

% matched lemmas: ambiente, apoio, crítico, futuro, habilidade, início, mencionar, pessoa, programa, quanto, que, questão, reconhecer, referência, ser
\end{textsample}
